\section{A weak realization of the linear subsidiary system}
\label{sec:constraint-existence}

This section concerns only the linear, first-order subsidiary system. It does
not prove existence for the geometric evolution or for its elliptic gauge.
The parameters are $a\kappa _1=4$, $\kappa _2=-99/100$, and the
Appendix-C scalar correction $\epsilon=1$. We use nondimensional $a=1$
when displaying its positive matrices. For general $a$, the value variables
in the energy are $\widehat\tau=\tau/a$ and
$\widehat\zeta=\zeta/a$; the wave variables already have the same dimensions.

Let $\mathcal B=\{r<S\}$ and let $\Omega=(S^2-r^2)/(2aS)$.
After the exact differential reduction, write the twenty-component equation as
\begin{equation}
 \partial_t U=\mathcal A U
 =V^i(x)\partial_iU+
 \frac{A(x)}{a\Omega(x)}J_0(n)U+R(x)U
 \qquad\hbox{in a collar of }\partial\mathcal B .
 \label{eq:subsidiary-abstract}
\end{equation}
Here $R$ is smooth and bounded up to the boundary. On the sixteen wave
components $V^i=B^iI+AP_i$, where the $P_i$ are the constant first-order
four-scalar-wave matrices; on the four value components $V^i=B^iI$.
The exact reduction, not merely an indicial calculation, supplies this
decomposition. The matrix $J_0$ is its boundary residue, with all bounded
differences absorbed into $R$.

The common metric is the identity on wave velocities and gradients, and on
the values has quadratic form
\begin{equation}
 4\widehat\tau^2+4\widehat\tau\widehat\zeta_n
       +8|\widehat\zeta|^2 .
 \label{eq:subsidiary-value-metric}
\end{equation}
Denote the resulting matrix by $H(n)$. The exact rational block calculation
establishes constants $h_*,h^*,\sigma>0$, independent of $n$, such that
\begin{equation}
 h_*I\le H\le h^*I,\qquad HV^i=(HV^i)^T,
 \qquad -\operatorname{sym}(HJ_0)\ge\sigma I .
 \label{eq:subsidiary-common-metric}
\end{equation}
The uniformity follows also from rotational covariance. No asserted
positivity of an uncomputed interior residue is needed.

Choose a smooth radial cutoff $\chi$, equal to one in a smaller boundary
collar and zero near the center. Extend $H$ smoothly and positively into the
interior, retaining the identity on the wave sector and a positive metric on
the value sector. This keeps $HV^i$ symmetric: value principal matrices are
scalar multiples of the identity. Rewrite the global operator as
\begin{equation}
 \mathcal A=V^i\partial_i+
       \chi\frac{A}{a\Omega}J_0+R_{\rm b},
 \label{eq:subsidiary-global-decomposition}
\end{equation}
where $R_{\rm b}$ is bounded. The cutoff transition is a positive distance
from scri, so this change introduces no additional singularity. One may
interpolate the value metric rotationally covariantly; the resulting metric
is smooth at the center and the global operator remains rotationally
covariant.

\paragraph{Proposition (linear weak subsidiary evolution).}
Assume the exact coefficients and positive matrix identities above. For every
$U_0\in L^2_H(\mathcal B)$ there is a unique
$U\in C([0,\infty);L^2_H)$ solving
\eqref{eq:subsidiary-abstract} distributionally on $\Omega>0$ with initial
value $U_0$. It obeys, on every bounded time interval,
\begin{equation}
 \|U(t)\|_H^2+
 c\int_0^t\!\int_{\mathcal B}
      \frac{\chi A}{a\Omega}|U|^2\,dx\,ds
 \le C_T\|U_0\|_H^2 .
 \label{eq:subsidiary-weak-energy}
\end{equation}
The solution operators form a strongly continuous semigroup. Its generator
is the maximal distributional realization
\begin{equation}
 D(\mathcal A)=\{U\in L^2_H:
          \mathcal A_{\rm formal}U\in L^2_H
          \text{ distributionally on }\Omega>0\}.
 \label{eq:subsidiary-maximal-domain}
\end{equation}
In particular, no independently chosen incoming boundary trace is added at
scri. This statement is about the first-order subsidiary extension; the
reduction-compatible subspace is invariant as proved below.

\paragraph{Proof: bounded approximants.}
Replace only the explicitly certified singular multiplier in
\eqref{eq:subsidiary-global-decomposition}, setting
\[
 q_N(\Omega)=\frac{1}{\sqrt{\Omega^2+N^{-2}}},\qquad
 \mathcal A_N=V^i\partial_i+
          \chi\frac Aa q_NJ_0+R_{\rm b}.
\]
For $\Omega>0$, $0\le q_N\le\Omega^{-1}$ and $q_N\uparrow\Omega^{-1}$.
Every approximant preserves precisely the same pointwise dissipativity
certificate. We do not cap an entire unverified-sign residue.

For each finite $N$, extend the smooth bounded coefficients and a positive
symmetrizer to a neighborhood of the closed ball, then to a bounded-coefficient
whole-space symmetric hyperbolic system. Ordinary Friedrichs regularization
and the bounded-coefficient energy estimate construct its $L^2$ Cauchy
solution. Restriction to the ball is independent of the exterior extension.
Indeed, at $r=S$ the wave normal coefficient $B_nI+AP_n$ has eigenvalues
$0,-2A$, while the value coefficient is $-A$. Thus its energy boundary flux
is nonpositive. A difference with zero data inside the ball has zero energy
there. The same argument first for smooth solutions and then by $L^2$
approximation justifies the restriction without prescribing extra boundary
data. Coefficients may be cut off farther outside the ball to ensure all
whole-space coefficients are bounded; this does not affect the restriction.

Integration by parts and \eqref{eq:subsidiary-common-metric} give
\begin{equation}
 \|U_N(t)\|_H^2+
 2\sigma\int_0^t\!\int\chi\frac Aa q_N|U_N|^2
 \le\|U_0\|_H^2+C\int_0^t\|U_N(s)\|_H^2\,ds .
 \label{eq:subsidiary-cap-energy}
\end{equation}
Here $C$ depends on bounded $R_{\rm b}$ and derivatives of $HV^i$, and is
independent of $N$. The potentially large derivatives of $q_N$ never enter
this zeroth-order energy estimate.

\paragraph{Proof: passage to the weak limit.}
On every finite time interval the approximants are bounded in
$L^\infty_tL^2_x$. On every compact subset of $\Omega>0$, their time
derivatives are bounded in $L^\infty_tH^{-1}_x$, uniformly in $N$.
A diagonal weak extraction therefore gives a limit, weakly continuous in
$L^2$ and solving the original equation distributionally in the open ball.
One may use a countable dense family of compactly supported smooth tests and
their uniformly equicontinuous pairings to obtain convergence at every time.
Density and the uniform $L^2$ bound give global weak continuity and the
correct weak initial trace. No strong convergence of a product of unknowns
is required: the equation is linear, and every coefficient is bounded on
each compact subset of the open ball.

For each fixed interior compact set, weak lower semicontinuity passes the
weighted part of \eqref{eq:subsidiary-cap-energy} to the limit. Exhaustion
by compact sets and monotone convergence then give the weighted integral in
\eqref{eq:subsidiary-weak-energy}. The same lower semicontinuity gives the
endpoint energy bound. This proves weak existence with the asserted energy.

\paragraph{Proof: weak uniqueness without assuming boundary traces.}
Let $\eta_\varepsilon(\Omega)$ vanish for $\Omega<\varepsilon$ and equal one
for $\Omega>2\varepsilon$, with
$0\le\eta_\varepsilon'\le C/\varepsilon$. Local Friedrichs mollification
justifies the energy identity with multiplier $\eta_\varepsilon H$ for an
$L^2$ weak solution: on its support all reaction coefficients are bounded,
the principal coefficients are Lipschitz, and their Friedrichs commutators
converge in $L^2$. This is a local statement; it does not differentiate the
singular multiplier across scri.

The only potentially positive cutoff error comes from the glancing normal
subspace. Indeed,
\[
 V^i\partial_i\Omega=-\gamma(B_nI+AP_n),\qquad
 \gamma=\frac r{aS},\qquad
 B_n+A=\frac{(S-r)^2}{2aS}=O(\Omega^2).
\]
Every other wave/value eigenvalue of $V^i\partial_i\Omega$ is nonnegative
in a sufficiently small collar. Hence the positive part of the cutoff error
is at most
\[
 C\varepsilon\int_{\{\varepsilon<\Omega<2\varepsilon\}}|U|^2\,dx.
\]
Its time integral tends to zero for $L^\infty_tL^2_x$ solutions. All negative
cutoff fluxes may be discarded. The certified reaction is dissipative and
its nonnegative integral passes by Fatou. Letting $\varepsilon\downarrow0$
therefore gives the same global energy inequality for a weak difference.
Zero initial data imply zero solution. No pointwise trace, selected sequence
of boundary spheres, or prior weighted trace estimate is required.

The weak initial continuity and upper energy bound imply strong continuity
at zero. Uniqueness and time independence give the semigroup property;
boundedness on finite time intervals then gives strong continuity at every
time. Thus the weak construction produces a $C_0$ semigroup, rather than
assuming one in advance.

\paragraph{Proof: identification of the maximal domain.}
For $U\in D_{\max}$ as defined in \eqref{eq:subsidiary-maximal-domain}, the
stationary local Friedrichs/cutoff argument gives
\begin{equation}
 \operatorname{Re}\langle\mathcal A_{\rm formal}U,U\rangle_H
 \le C\|U\|_H^2-
 \sigma\int\chi\frac A{a\Omega}|U|^2 .
 \label{eq:subsidiary-stationary-domain-energy}
\end{equation}
In particular the weighted integral is finite; it need not be imposed as a
separate domain assumption. The cutoff term has the same favorable sign
and $O(\varepsilon)$ positive bound as above.

The semigroup generator is contained in $D_{\max}$ by differentiating the
distributional equation against compactly supported tests. Conversely,
take $U\in D_{\max}$, choose real $\lambda>C$, and set
$F=(\lambda-\mathcal A_{\rm formal})U$. The Bochner integral
\[
 V=\int_0^\infty e^{-\lambda t}S(t)F\,dt
\]
belongs to the semigroup generator domain and satisfies
$(\lambda-\mathcal A_{\rm formal})V=F$. Therefore $W=U-V\in D_{\max}$
satisfies $(\lambda-\mathcal A_{\rm formal})W=0$.
Equation \eqref{eq:subsidiary-stationary-domain-energy} forces $W=0$.
This proves equality of the generator and maximal distributional domains.
It also makes the higher graph spaces precise: generator powers are the
corresponding recursively defined formal powers in $L^2$.

\paragraph{Reduction compatibility.}
Return to the unscaled value variables and define, distributionally in the
open ball,
\[
 E^Q_i=Q_i-2G_i\tau-\Omega\partial_i\tau
       =Q_i-\frac{\partial_iT}{\Omega},\qquad
 E^Z_{ij}=Z_{ij}-G_i\zeta_j-\Omega\partial_i\zeta_j
       =Z_{ij}-\partial_i z_j .
\]
The exact first-order equations, with no use of a Hamiltonian or momentum
projection, give
\begin{align}
 (\partial_t-B\cdot\partial)E^Q_i
  &=-\left[l+\frac{A(w+d_\theta)}\Omega\right]E^Q_i,\\
 (\partial_t-B\cdot\partial)E^Z_{ij}
  &=\left[\frac{Ac}\Omega-3l\right]E^Z_{ij},
 \qquad d_\theta=\kappa_1(2+\kappa_2),\quad c=2l-\kappa_1.
 \label{eq:subsidiary-reduction-transport}
\end{align}
All forcing cancels exactly. Backward characteristics of these equations
remain inside the ball: forward characteristics solve $\dot x=-B=x/a$.
On every compact portion of a backward characteristic the coefficients are
smooth and bounded. Local distributional uniqueness for transport thus
preserves zero initial reduction errors, without any data at scri.
The compatible subspace is closed in $L^2$ with its distributional
interpretation and invariant under $S(t)$. Consequently
$Z_{ij}=\partial_i z_j$ and
$\partial_jZ_{ij}=\partial_iZ_{jj}$ hold for compatible limits even though
the capped approximants need not preserve these identities.

\paragraph{The explicit joint phase and its interior generators.}
There is no need to assume that $\mathcal A_N^jU_0$ is bounded uniformly in
$N$; capping can break cancellations.  Maximal-domain identification instead
makes time-generator graph domains of the limiting semigroup invariant.
For the normal embedding we use a \emph{stronger} joint phase from the outset,
not just this time graph domain.  Set
\[
 \mathcal D=b(r)\partial_r,\qquad b=\Omega/\Omega_r
\]
on an outer collar, and cut it off smoothly away from the centre.  Let
$\chi\partial_i$ be three Cartesian generators supported away from scri,
with $\chi=1$ near the centre.  Choose overlapping cutoffs so that these
and $\mathcal D$, together with the three covariant rotations, span the
ordinary spatial derivatives uniformly on every compact interior region.
Write $\mathcal Z$ for this finite family.  At order $N$ the phase norm is
\[
 \mathcal G_N(U)=\sum_{p+|I|\le N}
  \|\mathcal Z^I\partial_t^pU\|_{L^2_H},
\]
where all words $I$ in the generators are included and time jets are the
actual limiting-generator jets.  Thus the phase explicitly includes
conormal and ordinary interior regularity.  For example $N=128$ controls
ordinary interior derivatives through order $128$.  It does not assert that
time/rotation graph regularity alone generates the conormal derivatives.

\paragraph{Propagation of the joint phase without a capped commutator.}
Write $\mathcal A=V^r\partial_r+V^A\partial_A+R$ in the collar, with
$R=A_*J_0(n)/(a\Omega)+R_b$.  The exact identity is
\[
 [\mathcal A,\mathcal D]=b'\mathcal A
 -(b'V^A+b\partial_rV^A)\partial_A
 -(\partial_rV^r)\mathcal D-b'R-b\partial_rR.
\]
Its apparent pole is removable, since $J_0(n)$ is radially constant and
\[
 b'\frac{A_*J_0}{a\Omega}
 +b\partial_r\frac{A_*J_0}{a\Omega}
 =\left[\frac{(b'-1)A_*}{a\Omega}
       +\frac{\partial_rA_*}{a\Omega_r}\right]J_0,
 \qquad b'-1=-\frac{\Omega\Omega_{rr}}{\Omega_r^2}.
\]
All remaining coefficients are smooth and bounded.  Replacing
$\mathcal AU$ by $U_t$ makes the conormal hierarchy triangular: a current
conormal derivative has the original generator plus bounded multiplication,
and its time/angular forcing has lower conormal order at the same total
joint order.  Rotational covariance gives the exact rotational commutations.
The interior Cartesian commutators have only bounded first-order
coefficients, supported where $\Omega$ is bounded below; expressing them in
the same overlapping generator family closes the ordinary interior jet
energy at the same total order.

Construct each conormal level using the limiting semigroup and bounded
perturbations, with its prescribed joint initial value.  On interior domains
the distributional derivative satisfies the same equation and initial data;
local uniqueness identifies it with the constructed $L^2$ field.  Backward
domains remain inside the ball at finite time, so no unknown derivative
trace at scri is required.  The finite augmented energy and completion in
$\mathcal G_N$ give an invariant strongly continuous joint phase.  This uses
no uniform commutator bound for the capped multipliers.

The previous unqualified claim that a time/rotation-only $\mathcal G_{40}$
supplied the six conormal derivatives is withdrawn.  The normal embedding
below uses the explicitly strengthened joint phase just defined.  It still
requires the fast/slow weight estimates and finite indicial-jet recursion;
no coupled geometric/gauge conclusion follows from the subsidiary
semigroup alone.

